\subsubsection{模型准备}\label{sec:q3-1}

\textbf{（1）沿用的记号与信息状态。}目标圆域 $\Omega$、频道 $f$、源位置 $p$、接收半径 $R_f$、机器狗当前位置 $x_t$、定位区域的外包多边形 $\widehat C_f$ 及其包围圆的圆心 $z(\widehat C_f)$、半径 $\rho(\widehat C_f)$ 都沿用第二问的记号。程序把 $20$ 个频道分成四类：还没收到过信号的\textbf{未发现}频道；收到过示向度、正在定位的\textbf{定位中}频道，每个保存一个外包多边形、它的包围圆、测过的位置，以及收到过信号和没收到信号的位置 $x^+$、$x^-$；收到真实成功反馈的\textbf{已清除}频道；以及已经能断定没有源的\textbf{无源}频道。任务开始时 $20$ 个频道全在未发现一类，任务结束时全部落在已清除或无源两类。源的真实位置、数量和接收半径都不进入决策，算法只通过检测和清除的反馈逐步了解它们。

\textbf{（2）总耗时与评价指标。}按题目规则，一局的总耗时由五项累加：
\begin{equation}\label{eq:q3-2}
T=\frac{L}{v}+5M+C+3A+2K,\qquad v=5\ \text{m/s},
\end{equation}
其中 $L$ 是总路程，$M$ 是检测次数，$C$ 是换频次数，$A$ 是光学定位的尝试次数（成功和失败都算），$K$ 是成功清除的次数。所以一次成功的光学清除共花 $5$ s，一次失败只花 $3$ s；清除动作不改变测向机当前的频道。题目的评价指标是平均定位清除时间 $T/N$，$N$ 为清除的源数；后面比较不同方案时，先对每个场景算 $T/N$，再对场景等权平均。

\textbf{（3）无信号反馈的用法。}第二问只用到有信号的反馈：示向度给出一个 $\pm1^\circ$ 的楔形，与原区域求交。第三问里机器狗在每个检测站都要把未发现频道扫一遍，对已发现的源也常常顺便再测一次，因此会积累不少“在某处对某频道没有收到信号”的记录。对全向源，无信号只有一个原因：离源超过了接收半径。这一点能从两方面缩小定位区域，见图\ref{fig:q3-negative}。

第一，接收半径至少 $1000$ m，所以以无信号点 $x^-$ 为圆心、半径 $1000$ m 的圆盘内不可能有源。把这个圆盘从定位区域里挖掉以后，剩下的部分一般不再是凸的，程序取它的凸包作为新的外包多边形，这样只会放大、不会缩小区域，真实源不会被裁掉。

第二，同一个源的接收半径在一局里固定不变。若在 $x^+$ 收到了信号、在 $x^-$ 没有收到，则源到 $x^+$ 的距离不超过 $R_f$，到 $x^-$ 的距离超过 $R_f$，于是源离 $x^+$ 比离 $x^-$ 更近：
\begin{equation}\label{eq:q3-3}
\|p-x^+\|\le R_f<\|p-x^-\|
\;\Longrightarrow\;
(x^--x^+)^{\mathsf T}p\le\frac{\|x^-\|^2-\|x^+\|^2}{2}.
\end{equation}
右边是一个半平面，边界正是 $x^+x^-$ 的垂直平分线。每来一个新的无信号点，就把它和该频道所有有信号点配对，逐个切掉半平面；程序在右端加了 $10^{-5}$ 的余量，把严格不等式放宽成闭半平面。这两条推断都只对全向源成立，第四问的定向源不能这样用。

\begin{center}
\begin{minipage}{\linewidth}
\centering
\QThreeNegativeFigure
\captionof{figure}{无信号反馈的两种裁剪。(a) 以无信号点为圆心的 $1000$ m 圆盘内不可能有源，挖掉后取凸包。(b) 源到有信号点的距离不超过接收半径、到无信号点的距离超过接收半径，所以源在两点垂直平分线靠近 $x^+$ 的一侧。}
\label{fig:q3-negative}
\end{minipage}
\end{center}

\textbf{（4）定位区域的构造与更新。}某个频道第一次在位置 $x$ 收到示向度 $\theta$ 时，定位区域取三者的交：以 $x$ 为顶点、沿 $\theta$ 方向张开 $\pm\varepsilon$、长 $1500$ m 的三角形（它包含半径 $1500$ m 的扇形），场地的外切正 $64$ 边形，以及该频道此前所有无信号点按（3）给出的裁剪。以后每收到一次示向度，就再交上一个楔形，并重新做一遍无信号裁剪。程序里 $\varepsilon$ 取 $1.005^\circ$，比题面的 $1^\circ$ 多出的 $0.005^\circ$ 用来容纳示向度保留两位小数带来的舍入。每次更新后重新计算包围圆，方法是枚举多边形顶点中由一个、两个或三个顶点支撑的圆，取半径最小且包住全部顶点的那个；区域顶点很少，枚举很快。

\textbf{（5）不能省的三条底线。}后面所有的规则都在这三条之内做取舍：第一，真实源始终在它的定位区域里，区域只按真实反馈和上面两条推断收缩；第二，一个未发现频道只有在“它在场地里任何位置都会被某个已扫描的站收到”的条件下，才能判为无源；第三，只有收到设备返回的清除成功，才把频道记为已清除。程序把这三条写成运行时检查，任何一条被违反就报错停机，而不是继续跑下去。
